\chapter{Den siste metoden}

\begin{motto}
Stilklassifisering var det næraste designhistoria kom kvantitet.\\ Ei maskin gjer det no på brøkdelen av eit sekund. Spørsmålet er kva faget då har att.
\end{motto}

Eit fag kan leve lenge utan ein teori om det berre eig ein \emph{metode}: ein etterprøvbar prosedyre som skil rett frå gale og lèt to forskarar komme til same svar. Designhistoria har, gjennom heile si akademiske levetid, hatt nett éin. Å plassere eit objekt i ein periode, attribuere det til ein meister, datere det, byggje typologiar av slektskap; dette var det næraste faget kom noko som kjendest som kunnskap og ikkje som smak. Kapitlet handlar om kva som hende då ei maskin lærte å gjere det same. Stilklassifiseringa vart ikkje motbevist. Ho vart \emph{automatisert}, og då ho var det, viste det seg kor lite ho hadde vore: ikkje katalogen over ei djup innsikt, men mønstergjenkjenning på overflatestatistikk, slik ein ornitolog kjenner art på fjørdrakt. Linné, ikkje Darwin. Designhistoria har stogga ved Linné i hundre år og kalla det å forstå form \autocite{fallan2010}.

\section{Kva stilklassifiseringa eigentleg var}

Sjå nøkternt på operasjonen. Stilhistoria tek eit korpus av objekt og sorterer dei etter visuelle trekk: linjeføring, ornament, materiale, proporsjon. Ut av sorteringa kjem periodar, attribusjonar (dette er Aalto, dette er Eames) og dateringar. Dette er ekte arbeid og krev eit trent auge. Men spør kva slag operasjon det er, reint formelt, og svaret er ubehageleg. Stilen til eit objekt er korrelasjonsstrukturen i trekka, kva som plar opptre saman; å kjenne att ein stil er å kjenne att eit statistisk mønster. Så lenge berre menneske kunne gjere det, slapp faget å spørje. Det trente auget verka magisk, kjennaren såg det novisen ikkje såg, og evna lét seg ikkje setje ord på. Men at ei evne er vanskeleg og taus, tyder ikkje at ho er djup. Det kan like gjerne tyde at ho er reint statistisk, ei evne hjernen byggjer av tusen døme utan å gå vegen om eit omgrep. Og den slags evne gjer maskinene betre enn oss.

I dag er stilklassifisering eit løyst problem i datasynet. Eit konvolusjonelt nettverk, trena på eit biletkorpus, klassifiserer stil, attribuerer opphav og daterer objekt med ein presisjon som matchar fagfolk, på brøkdelen av eit sekund \autocite{saleh2016classification}. Leon Gatys og kollegaene hans synte endå skarpare kva stil \emph{er} i denne samanhengen: dei skilde stil frå innhald i eit bilete reint matematisk, og fann at stilen lèt seg fange fullstendig i feature-statistikken til nettverket, i korrelasjonane mellom aktiveringar samla i ein Gram-matrise \autocite{gatys2016style}. Stil, teknisk, er eit sett gjennomsnitt og samvariasjonar. Difor lèt han seg overføre frå eitt bilete til eit anna med ein algoritme. Og difor kan ei maskin lese kjennaren sitt trente auge som det det var: ein langsam, dyr og udokumentert versjon av ein klassifikator det no tek minutt å trene.

\section{Handa til van Gogh}

Kor lite stil var, syner seg i kor lett maskina flyttar han. I 2015 gav Gatys og kollegaene hans nettet eit vanleg foto, ei gate, og ved sida av det «Stjernenatt», van Goghs himmel frå 1889 med dei virvlande, tjukke stroka. Så bad dei algoritmen om noko som høyrest umogleg ut: hald innhaldet i fotoet, men mal det i van Goghs stil. Maskina gjorde det på sekund. Ut kom gata, framleis attkjenneleg, men no med ein himmel som krøllar seg og strok som lyser, med heile den umiskjennelege handa til ein daud nederlendar lagd over ein augneblink han aldri såg. Innhaldet, kvar kantane og objekta sit, ligg i dei høgare laga av nettet. Stilen er noko anna og merkeleg fattig: samvariasjonen i aktiveringane, kva strok og fargar og teksturar som plar opptre saman, og kor mykje. Difor let han seg skilje frå motivet og leggje over eit anna bilete. Stil er separabel fordi stil er statistikk, og statistikk kan lyftast ut og flyttast.

No held dette opp eit spegelbilete av designhistorikaren på sitt mest vitskaplege. Han ser eit ukjent objekt, ein stol, og plasserer han: seint femtital, skandinavisk, i slekt med Wegner. Han kjenner att eit mønster av trekk som plar opptre saman, beinform, materialval, fugemåte, profil. Det er Gram-matrisa lesen av eit menneske i staden for ei maskin, ei kjensle for kva som høyrer saman, bygd av tusen døme utan eit omgrep. Stiloverføringa er ikkje ein kuriositet frå nabofaget. Ho syner at det som kjendest som den djupaste forma for å forstå form, var den lettaste å rekne ut.

\begin{kasus}{Maskina som las kunsthistoria utan å få lese ho}
I 2018 sette eit forskarteam eit nettverk til å ordne tusenvis av verk frå den vestlege målerihistoria, frå tidleg renessanse til samtid, utan å gje det periodenamn, lærebok eller noko omgrep om kva ein stil er \autocite{elgammal2018shape}. Det fekk berre bileta. Då det var ferdig, hadde det lagt dei langs ein akse som nær fullkome følgde kronologien, i same rekkjefølgje ein kunsthistorikar ville gjeve etter eit langt studium. Maskina rekonstruerte grovstrukturen i stilhistoria frå rein pikselstatistikk, utan å skjøne eit einaste måleri. Det kunsthistorikaren har halde for ein djup, trent dom, fall ut av nettet som biprodukt av å sjå på pikslar.
\end{kasus}

\section{Stil mot form}

Her ligg det skiljet kapitlet eigentleg dreier om, og det er ikkje skiljet mellom menneske og maskin. Det er skiljet mellom to ting ein kan gjere med eit objekt. Ein kan klassifisere \emph{stilen} hans, overflatestatistikken, det maskina alt eig. Eller ein kan måle \emph{forma} hans, den fysiske geometrien, og det er noko anna slag.

For form lèt seg gjere til eit kvantitativt objekt, på sine eigne premiss, og det finst eit gammalt døme på korleis. I staden for å sortere etter stil kan ein definere eit \emph{formrom} der kvar akse er ein målbar eigenskap, og leggje kvart objekt inn som eit punkt med koordinatar. Dette er morfospacet, eit samanhengande landskap forma kan gli gjennom, ikkje den diskrete katalogen formgrammatikken spytta ut i kapittel tolv, sjølv om same nemninga dekkjer båe. Då vert variasjon til geometri, slektskap til avstand, endring over tid til ei bane ein kan rekne på. D'Arcy Thompson skisserte tanken for den levande forma alt i 1917 \autocite{thompson1917growth}; David Raup la i 1966 heile mangfaldet av skjelhus inn i eit rom med tre aksar og kunne rekne på det \autocite{raup1966}; evolusjonsbiologien bygde tilpassingslandskapet sitt oppå same grepet \autocite{wright1932}. Dette er ikkje malen design skulle ha nått. Det er berre eit prov på at form \emph{kan} målast, og at fråvêret av ei slik måling i design ikkje gjer faget til ein mislukka biologi, men til eit fag som valde å gjere noko anna med objekta sine.

Skilnaden vert tydeleg i det fysiske. Sjå eit lager fullt av stolar, tusenvis, frå to hundre år. Designhistorikaren går gjennom dei med kataloglogikken: denne er jugend, denne Bauhaus, denne midtårhundre-moderne. Kvar stol får eit kort med ein periode og ei slekt, ein taksonomi av etikettar, og det er, vi veit det no, nett operasjonen ei maskin gjer betre. Morfometrikaren rører ikkje etikettane. Han måler. Setehøgd i millimeter, ryggvinkel i grader, forholdet mellom setedjupn og setebreidd, krumninga i ryggbøylen som ein parameter, talet på berande element. Han legg kvar stol inn som eit punkt i formrommet. No er stolane ikkje ei samling etikettar lenger, men ei sky av punkt, og skya har ein struktur.

Og her kjem det funnet som søsterverket til denne boka demonstrerer empirisk, og det er kapitlet sitt skarpaste. Legg fire tusen stolar inn i eit slikt formrom og spør om stiletiketten føreseier kvar punktet hamnar. Han gjer det berre svakt: stilnamnet forklarer mindre enn ein femdel av variasjonen langs hovudaksen, den som spenner frå låg og brei til høg og smal. Stolar med same stilnamn ligg spreidde; stolar frå ulike periodar ligg side om side når geometrien deira liknar. Den verkelege variasjonen, den som styrer korleis ein stol ber, kjennest og held, ligg i den målbare forma og ikkje i etiketten. Stilen, som faget bygde heile si kvantitative sjølvaktelse på, fangar mindre av kva ein stol faktisk \emph{er} enn tre tal tekne med eit måleband.

Designhistoria tok aldri det steget. Ho heldt fram med å klassifisere stil, overflata, og difor kan maskinene no overta arbeidet hennar, for dei er meister på overflate. Det er eit val om kva slag fag design skulle vere. Faget valde aksen der maskina alt rår, og lét liggje aksen der maskina kunne vore tenaren, der ein algoritme målte fire tusen stolar over natta og gav historikaren eit kvantitativt prosjekt han eigde. Skiljet mellom stil og form er skiljet mellom det maskina alt har teke og det faget kunne ha eigd.

\section{Kva som no står att}

Når stilklassifiseringa fell bort som faget sitt eige bidrag, kva er att? Dei andre utvegane var prøvde og forkasta gjennom heile denne delen: forskinga som ikkje akkumulerer, atelierdommen som er smak, dei lånte metodane som er andre fag sine. Att stod stilklassifiseringa, det eine grepet som var faget sitt eige og som likna ein metode. No er ho maskina sin. Eit fag i den stoda kan framleis produsere tekstar, konferansar og grader; det kan ikkje produsere kunnskap som hopar seg opp og let seg prøve.

\vignett

Faget har altså mist sin siste eigne metode, og den utvegen som fanst, formmålinga, krev nett det fundamentet faget aldri bygde. Men kanskje finst det forsvararar som vil seie at sjølve kravet om ein kvantitativ metode er urettferdig, at design eig andre, like gyldige former for kunnskap. Det er den sterkaste innvendinga mot heile denne boka, og ho fortener eit eige oppgjer.

\laeringsmal{\begin{itemize}
\item Forklare kvifor stilklassifisering var faget sin einaste kvantitative metode, og kvifor operasjonen i botnen berre er mønstergjenkjenning på overflatestatistikk.
\item Gjere greie for at «stil» teknisk er feature-statistikk (Gatys), og difor eit løyst maskinlæringsproblem.
\item Skilje stil (overflatestatistikk) frå form (målbar geometri), og forklare kvifor faget valde katalogen framfor formrommet.
\end{itemize}}

\ovingar{\begin{enumerate}
\item Skisser korleis du ville bygd ein stilklassifikator på eit designkorpus (t.d. stolar eller plakatar). Kva «veit» han når han er ferdig, og kva veit han ikkje?
\item Vel ti objekt i éin klasse og definer tre \emph{målbare} aksar (ikkje stiletikettar). Plott objekta. Kva ser du som stilhistoria ikkje fangar?
\item Om stilklassifisering er automatisert vekk, kva legitimt kvantitativt forskingsprosjekt står att for designhistoria? Grunngje svaret.
\end{enumerate}}

\vidarelesing{\fullcite{saleh2016classification}; \fullcite{gatys2016style}; \fullcite{elgammal2018shape}; \fullcite{raup1966}; \fullcite{thompson1917growth}; \fullcite{wright1932}; \fullcite{fallan2010}.}
